% Relatório de estudo sobre a otimização (02/11/2026, pedido do autor).
% Material de estudo, não é parte do TG. Todo número vem de
% dados/resultados/otimizacao_teste_2026-11-02/ (ferramentas/scripts/otimizacao_teste.sh),
% que lê os pareados da refação de 05/10/2026 (dados/resultados/refeito_2026-10-05/pareados/).
% Parte 1 de 3: otimização x teste. Mapa de calor e queda por forma de uso
% ficam para depois de o autor ler esta (regra nº 3).
\documentclass[11pt,a4paper]{article}
\usepackage[T1]{fontenc}
\usepackage[brazilian]{babel}
\usepackage{amsmath,amssymb}
\usepackage{mathptmx}
\usepackage[a4paper,margin=2.4cm]{geometry}
\usepackage{booktabs,array}
\usepackage[font=small,labelfont=bf]{caption}
\usepackage[most]{tcolorbox}
\usepackage{pgfplots}
\usepgfplotslibrary{groupplots}
\usetikzlibrary{shapes.geometric}
\pgfplotsset{compat=1.18}
\usepackage[hidelinks]{hyperref}
\setlength{\parskip}{0.45em}
\setlength{\parindent}{0pt}

\definecolor{cgarch}{RGB}{120,120,120}
\definecolor{carima}{RGB}{230,140,0}
\definecolor{cesc}{RGB}{220,50,32}
\definecolor{cmod}{RGB}{0,90,181}

\newcommand{\arq}[1]{../../../dados/resultados/otimizacao_teste_2026-11-02/#1}
\input{\arq{valores.tex}}
% número com vírgula decimal (os macros vêm com ponto)
\newcommand{\num}[1]{\pgfmathprintnumber[fixed,zerofill,precision=2,use comma]{#1}}

\pgfplotsset{base/.style={table/col sep=comma, grid=major, grid style={gray!20},
    tick label style={font=\scriptsize}, label style={font=\footnotesize},
    title style={font=\footnotesize\bfseries, align=center},
    /pgf/number format/use comma, unbounded coords=discard,
    width=7.4cm, height=5.6cm}}

\newtcolorbox{ideia}[1][]{colback=cmod!5, colframe=cmod!70, fonttitle=\bfseries, title={#1}, arc=2pt, boxrule=0.6pt}
\newtcolorbox{cuidado}[1][]{colback=cesc!4, colframe=cesc!70, fonttitle=\bfseries, title={#1}, arc=2pt, boxrule=0.6pt}
\newtcolorbox{banca}{colback=black!3, colframe=black!50, fonttitle=\bfseries, title=Se a banca perguntar, arc=2pt, boxrule=0.6pt}

% \painel{arquivo}{métrica: sharpe|calmar}{x da escolhida}{y da escolhida}{mediana do teste}{título}
\newcommand{\painel}[6]{%
  \nextgroupplot[title={#6}]
  \addplot[only marks, mark=*, mark size=1.1pt, cgarch, fill opacity=0.45, draw opacity=0]
    table[x=#2_otim, y expr={ifthenelse(\thisrow{garch}==1,\thisrow{#2_teste},nan)}]{\arq{#1}};
  \addplot[only marks, mark=square*, mark size=1.8pt, carima]
    table[x=#2_otim, y expr={ifthenelse(\thisrow{garch}==0,\thisrow{#2_teste},nan)}]{\arq{#1}};
  \draw[cmod, dashed, thick] ({rel axis cs:0,0}|-{axis cs:0,#5}) -- ({rel axis cs:1,0}|-{axis cs:0,#5});
  \addplot[only marks, mark=star, mark size=4.5pt, cesc, very thick] coordinates {(#3,#4)};
}

\title{\textbf{A otimização prevê o teste?}\\[2pt]\large as 312 configurações da PETR4, na otimização e no teste}
\author{Material de estudo para os autores do TG (não faz parte do texto)}
\date{2 de novembro de 2026 \quad$\cdot$\quad parte 1 de 3}

\begin{document}
\maketitle

\begin{ideia}[Em uma página]
\textbf{A pergunta.} A configuração escolhida como a melhor na otimização continua entre as melhores no período de
teste? Se continuar, a otimização serve para escolher. Se não, a escolha foi em boa parte sorte.

\textbf{Como se responde.} No período de teste, o robô só executou a configuração escolhida. Mas na refação de
05/10/2026 a grade inteira também foi rodada no teste, então dá para ver como teriam ido \emph{todas} as 312
configurações e onde a escolhida ficou entre elas.

\textbf{A resposta curta.} Pelo Sharpe, a otimização não prevê o teste: nos três períodos, a escolhida ficou
abaixo da mediana no teste. Pelo Calmar, prevê um pouco: a escolhida ficou acima da mediana em dois dos três.
\end{ideia}

\section{O que o gráfico mostra}

O MetaTrader tem uma aba chamada \emph{Gráfico de otimização}, que desenha o resultado de cada passagem da grade.
A Figura~\ref{fig:disp} usa a mesma ideia, com uma diferença: cada configuração aparece com \emph{dois}
resultados, um em cada eixo.

\begin{itemize}\setlength\itemsep{2pt}
  \item \textbf{Cada ponto é uma configuração} da grade do TG (só compra, sem o intervalo de confiança):
        \nA{} no total, \ngarchA{} com GARCH (cinza) e 8 sem GARCH (laranja).
  \item \textbf{Eixo horizontal:} o resultado nos três anos de otimização, que foi o que se usou para escolher.
  \item \textbf{Eixo vertical:} o resultado no período de teste, que a escolha não viu.
  \item \textbf{Estrela vermelha:} a configuração escolhida (o ARIMA-GARCH Sharpe na coluna da esquerda, o
        ARIMA-GARCH Calmar na da direita). Por definição, ela é o ponto mais à direita entre os cinza.
  \item \textbf{Linha azul tracejada:} a mediana do teste. Metade das configurações ficou acima dela e metade
        abaixo. É o que se esperaria, em média, escolhendo uma configuração ao acaso.
\end{itemize}

\textbf{Como ler.} Se a otimização previsse o teste, a nuvem subiria da esquerda para a direita, como uma
diagonal, e a estrela ficaria lá no alto. Se a nuvem não tem inclinação, ir bem na otimização não diz nada sobre
o teste.

\begin{figure}[p]
\centering
\begin{tikzpicture}
\begin{groupplot}[base, group style={group size=2 by 3, horizontal sep=1.6cm, vertical sep=2.7cm},
    xlabel={na otimização}, ylabel={no teste}]
  \painel{pontos_T1.csv}{sharpe}{\gsxA}{\gsyA}{\medSA}{Sharpe\\otimização 2004--2006 $\to$ teste 2007--2011}
  \painel{pontos_T1.csv}{calmar}{\gcxA}{\gcyA}{\medCA}{Calmar\\otimização 2004--2006 $\to$ teste 2007--2011}
  \painel{pontos_T2.csv}{sharpe}{\gsxB}{\gsyB}{\medSB}{Sharpe\\otimização 2010--2012 $\to$ teste 2013--2016}
  \painel{pontos_T2.csv}{calmar}{\gcxB}{\gcyB}{\medCB}{Calmar\\otimização 2010--2012 $\to$ teste 2013--2016}
  \painel{pontos_T3.csv}{sharpe}{\gsxC}{\gsyC}{\medSC}{Sharpe\\otimização 2022--2024 $\to$ teste 2025--set.~2026}
  \painel{pontos_T3.csv}{calmar}{\gcxC}{\gcyC}{\medCC}{Calmar\\otimização 2022--2024 $\to$ teste 2025--set.~2026}
\end{groupplot}
\end{tikzpicture}

\vspace{4pt}
{\small\sloppy
\tikz\fill[cgarch, opacity=0.6] (0,0) circle (1.6pt); com GARCH \quad
\tikz\fill[carima] (0,0) rectangle (4pt,4pt); sem GARCH \quad
\tikz\node[cesc, star, star points=5, fill, inner sep=1.6pt]{}; escolhida \quad
\tikz\draw[cmod, dashed, thick] (0,0) -- (0.6,0); mediana do teste}

\caption{Cada configuração na otimização (horizontal) e no teste (vertical), PETR4. Sharpe e Calmar calculados
com custos e caixa parado na Selic, como na escolha.}
\label{fig:disp}
\end{figure}

\section{O que deu}

A Tabela~\ref{tab:pos} resume os seis painéis. A \emph{posição} é o lugar da escolhida quando as \nA{}
configurações são ordenadas pelo resultado no teste: 1\textordfeminine{} é a melhor, \nA\textordfeminine{} é a pior.

\begin{table}[ht]
\centering
\caption{Onde a escolhida ficou no teste, entre \nA{} configurações}
\label{tab:pos}
\small
\begin{tabular}{l rr r rr r}
\toprule
 & \multicolumn{3}{c}{Escolha pelo Sharpe} & \multicolumn{3}{c}{Escolha pelo Calmar} \\
\cmidrule(lr){2-4}\cmidrule(lr){5-7}
Teste & posição & acima da mediana? & $\rho$ & posição & acima da mediana? & $\rho$ \\
\midrule
2007--2011       & \posSA\textordfeminine{} & não & \num{\rhoSA} & \posCA\textordfeminine{} & sim & \num{\rhoCA} \\
2013--2016       & \posSB\textordfeminine{} & não & \num{\rhoSB} & \posCB\textordfeminine{} & sim & \num{\rhoCB} \\
2025--set.~2026  & \posSC\textordfeminine{} & não & \num{\rhoSC} & \posCC\textordfeminine{} & não & \num{\rhoCC} \\
\bottomrule
\end{tabular}

\vspace{2pt}
{\footnotesize $\rho$: correlação de Spearman entre a ordem na otimização e a ordem no teste, nas \nA{} configurações.}
\end{table}

\textbf{O que é o $\rho$ (correlação de Spearman).} Ordenam-se as configurações duas vezes, pela otimização e pelo
teste, e mede-se o quanto as duas listas concordam. Vale 1 se a ordem é a mesma, 0 se não há relação nenhuma e
$-1$ se uma é o inverso da outra. Ele resume numa conta a inclinação da nuvem.

\textbf{Pelo Sharpe}, $\rho$ fica entre \num{\rhoSB} e \num{\rhoSA}: nenhuma relação. A escolhida ficou abaixo da
mediana nos três testes. No mais recente, ficou em \posSC\textordfeminine{} de \nA{}, quase a última. Isso não quer dizer que o
Sharpe seja um critério ruim em si. Quer dizer que, nesta grade e neste ativo, o Sharpe de três anos não carrega
informação sobre o Sharpe dos anos seguintes.

\textbf{Pelo Calmar}, $\rho$ fica entre \num{\rhoCB} e \num{\rhoCA}: uma relação fraca, mas no sentido certo. A
escolhida ficou acima da mediana em 2007--2011 e em 2013--2016 (\posCB\textordfeminine{} de \nA{}), e abaixo em 2025--2026. O Calmar
mede ganho por unidade de queda máxima, e a queda depende muito de \emph{como} o robô reduz a posição, que é uma
característica da configuração e não do acaso do período. Essa é uma explicação provável para ele se manter mais,
não algo que o gráfico prove.

Os pontos laranja (sem GARCH) não concorrem à escolha do ARIMA-GARCH, que é feita só entre os cinza. Estão no
gráfico para comparação. Com só 8 pontos, não dá para dizer mais nada sobre eles aqui.

\begin{banca}
\textbf{``O resultado de vocês não é só a sorte de ter escolhido a melhor configuração?''}

É exatamente por isso que o teste foi feito num período que a escolha não viu. E este gráfico mostra que a
preocupação tem fundamento: escolhida pelo Sharpe, a melhor da otimização ficou abaixo da mediana nos três testes.
Por isso o Cap.~III mede o desempenho em períodos posteriores à escolha, com a configuração congelada. A única
exceção é 2005--2026, e o texto já avisa que esse teste inclui os anos da escolha.

\textbf{``Então por que usar o Sharpe?''} Porque é a medida padrão de ganho por risco, e o critério foi
definido antes de ver os resultados. O gráfico mostra que o Calmar, que olha para a queda máxima, se mantém mais
entre períodos. Isso combina com o foco do trabalho em risco: a parte que a otimização consegue prever é a
proteção, não o retorno.
\end{banca}

\begin{cuidado}[O que este gráfico não diz]
\begin{itemize}\setlength\itemsep{1pt}
  \item São \textbf{três pares de períodos de um ativo só}. O padrão aparece nos três, mas três é pouco.
  \item Sharpe e Calmar aqui são os da \textbf{escolha} (com custos e caixa na Selic). O Cap.~III mostra o
        resultado bruto, então os valores do eixo não são os números das tabelas do TG.
  \item O período 2014--2016 $\to$ 2017--2024 (o teste principal) \textbf{não está aqui}, porque a grade não foi
        rodada inteira em 2017--2024. Só os robôs escolhidos e as referências foram.
  \item Em 2007--2011 a escolha pelo Sharpe e pelo Calmar deu \textbf{a mesma configuração}, por isso as duas
        estrelas de cima representam o mesmo robô.
  \item O TG hoje não fala de persistência entre otimização e teste. Se o autor quiser levar isto ao texto, é
        decisão dele, e as frases saem daqui só depois de ele dominar o gráfico.
\end{itemize}
\end{cuidado}

\section*{De onde vêm os números}

{\small\sloppy
Script \texttt{ferramentas/scripts/otimizacao\_teste.sh}, que lê
\texttt{dados/resultados/refeito\_2026-10-05/pareados/PETR4\_C\textit{n}\_T\textit{n}.csv} e grava em
\texttt{dados/resultados/otimizacao\_teste\_2026-11-02/}. Conferências: (1) a escolha refeita pelo script
coincide com a oficial (\texttt{escolhas/r3\_C\textit{n}.txt}) nos três períodos; (2) o Sharpe e o Calmar da
escolhida no teste coincidem com os do teste individual dela (\texttt{testes/metricas\_R3\_T\textit{n}\_G*.csv});
(3) com as 32 configurações do intervalo de confiança de volta (344), o $\rho$ reproduz o bloco 5 de
\texttt{analise\_grades.txt} (0,10; $-0{,}05$; $-0{,}02$ e 0,46; 0,28; 0,37).
\par}

\end{document}
